% !PRINTED 478
\weberpage{478}
%<<CONT
minante $-4m$ gehöriger quadratischer Faktor abzusondern, und
wenn $m=1$ ist, so ist die nach Absonderung der fremden
Faktoren übrig bleibende Gleichung ein Quadrat.

Als Beispiele für diese Fälle nehmen wir:

1. \quad $m=3,\ n=3,$ \quad Absonderung eines Linearfaktors.\\
2. \quad $m=2,\ n=5,$ \quad kein fremder Faktor.\\
3. \quad $m=2,\ n=3,$ \quad Absonderung eines quadratischen Faktors.\\
4. \quad $m=1,\ n=3,$ \quad Quadrat.\\
5. \quad $m=1,\ n=5,$ \quad Quadrat nach Absonderung eines quadratischen Faktors.\\
6. \quad $m=1,\ n=7,$ \quad Quadrat.

Im Falle 1. hat man in der auf den Transformationsgrad $n=3$
bezüglichen Formel des § 73 zu setzen:
\[
u^3=2,\quad v^3=f(\sqrt{-27})^3=2x,
\]
also
\[
A=x^2+\frac{1}{x^2},\quad B=4x-\frac{2}{x},
\]
und folglich:
\[
x^4-4x^3+2x+1=(x-1)(x^3-3x^2-3x-1)=0,
\]
so daß man für $m=27$ erhält:
\weq{1}{f(\sqrt{-27})^3=2x,\quad x^3-3x^2-3x-1=0.}

Im Falle 2. setzen wir im System II. des § 73 ($n=5$):
\[
u_1=f_1(\sqrt{-2})=\sqrt[4]{2},\quad v_1=f_1(\sqrt{-50})=\sqrt[4]{2}\,x,
\]
\[
A_1=\frac{1}{x^3}-x^3,\quad B_1=2\Bigl(x^2+\frac{1}{x^2}\Bigr),
\]
so daß man (ohne fremden Teiler) die Gleichung 6ten Grades:
\weq{2}{\frac{1}{x^3}-x^3+2\Bigl(x^2+\frac{1}{x^2}\Bigr)=0,\quad f_1(\sqrt{-50})=\sqrt[4]{2}\,x}
erhält, die sich für
\[
y=\frac{1}{x}-x
\]
auf den dritten Grad reduziert:
\weq{3}{y^3+2y^2+3y+4=0.}
